\section{Problem Formulation}\label{sec2}
%Before introducing our proposed models, we first formally detail some concepts 
%related to our work.

%Knowledge graph reasoning is the core of making explainable 
%recommendation[PGPR]. 
%In this paper, we consider to add temporal information to the inferred
%multi-hop paths on the KG. Accordingly, the knowledge graph reasoning for 
%explainable recommendation problem can be formulated as follows:
%
%\textbf{Input.} The model input consists of the user set $U$, the item set 
%$V$,

\begin{table*}[t!] 
    \centering
    \caption{Notation and Definitions.}
\begin{tabular}{cc} 
\toprule
\textbf{Notations} & \textbf{Annotation} \\ \hline
$\mathcal{U}$                  & User set                \\
$\mathcal{V}$                  & Global item set                \\ 
$\mathcal{V}_\mathcal{U}$      & The observed interactions between users and items  \\
$\mathcal{V}_\mathcal{U}^T$     & The clustered time-aware interaction set, with $L$ categories \\
$\mathcal{G}_T$                 & Time-aware Collaborative Knowledge Graph \\
$\mathcal{E}^{\prime}$          & The entities in Time-aware Collaborative Knowledge Graph \\
$\mathcal{R}_T$                 & The relations in Time-aware Collaborative Knowledge Graph \\
$T$                             & The interaction timestamp set \\
$\hat{T}$                       & The recommend time \\
$\hat{\mathcal{V}}$             & The recommended item set \\
$\tau_{u, \hat{\mathcal{V}}}$   & The reasoning path set for user $u$ to  item set $\hat{\mathcal{V}}$ \\
$\mathcal{T}$                & Temporal feature space, including statistical features and structural features \\
$r_{kg}$                        & The  external knowledge graph relation set \\
$r_{interact}$                  & The time-aware interaction relation set \\
$\Mat{r}$                       & The embedding of relations            \\
$\Mat{e}_h, \Mat{e}_t$          & The embedding of entities                \\
$\Mat{W}_{t},  \Mat{W}_{\hat{T}}$   & The weight distribution of the Gaussian models at timestamp $t$,\ $\hat{T}$ \\
$\Mat{W}_{h_u}$                 & The time cluster weight distribution of user $u$'s interaction history \\
$k, k^{\prime}, K$                  & The $k$-th reasoning step, state history length $k^{\prime}$ and terminal step \\
$\Mat{s}_k$                     & The state at step $k$ \\
$a_k, \Set{A}_k$                     & The action and action space at step $k$ \\
$g_R(\hat{v} \mid u)$           &  The time-aware based reward for recommended item $\hat{v}$ for user $u$ \\
$\pi\left(\cdot \mid \mathbf{s}_k, \mathcal{A}_{k}(u)\right)$ & The policy network at step $k$ \\
$\hat{c}(\mathbf{s}_k)$         & The value network at step $k$ \\ \bottomrule
\end{tabular}
\label{notation_table}
\end{table*}


\noindent\textbf{Notions:} 
We consider the time-aware path reasoning for KG-enhanced recommendation.
Let denote the user set, the item set, and the observed interactions as $\mathcal{U}$, $\mathcal{V}$, ${\mathcal{V}}_{\mathcal{U}}$, respectively.
Moreover, there exists an external KG storing item properties. Generally, the KG is represented as $\mathcal{G}=\{(h, r, t)  \mid h, t \in \mathcal{E}, r \in \mathcal{R}\}$, where $\mathcal{E}$ represents the set of real-world entities and $\mathcal{R}$ is the set of relation.  
Each triplet $(h,r,t)$ delineates that there exist a relationship $r$ from head entity $h$ to tail entity $t$.
Previous works \cite{multi_modal,kgat} combine the user behaviors and item knowledge as a unified relational graph, termed Collaborative Knowledge Graph (CKG), which is defined as 
$\mathcal{G}_R=\{(h, r, t) \mid h, t \in \mathcal{E}^{\prime}$, $ r \in \mathcal{R}^{\prime}\}$,  where $ \mathcal{E}^{\prime}=\mathcal{E} \cup \mathcal{U}$ and $ \mathcal{R}^{\prime}=\mathcal{R} \cup{\mathcal{V}}_{\mathcal{U}}$. The notations of this work are summarized in Table \ref{notation_table}.
	 
\vspace{5pt}
\noindent\textbf{Reasoning Environment:}
Considering the time-aware interactions instead, we further obtain a Time-aware Collaborative Knowledge Graph (TCKG), which is the environment of time-aware path reasoning.
Specifically, the interaction timestamps are grouped into $L$ groups after time-series clustering, and the interaction relation ${\mathcal{V}}_{\mathcal{U}}$ is extended to $\{\mathcal{V}_{\mathcal{U}}^1, 
\mathcal{V}_{\mathcal{U}}^2, \cdots, \mathcal{V}_{\mathcal{U}}^{L}\}$.
Hence, the TCKG is formulated as $\mathcal{G}_T=\{(h, r, t) \mid 
h, t \in \mathcal{E}^{\prime}, r \in \mathcal{R}_{T}\}$, where 
$ \mathcal{R}_{T}=\mathcal{R} \cup \{ \mathcal{V}_{\mathcal{U}}^1, 
\mathcal{V}_{\mathcal{U}}^2, \cdots, \mathcal{V}_{\mathcal{U}}^{L} \}$.


\vspace{5pt}
\noindent\textbf{Task Description:}
We present the task of time-aware path reasoning for recommendation (TPRec) as follows:
\begin{itemize}
\item \textbf{Input:} Time-aware collaborative knowledge graph $\mathcal{G}_T$;
% \item \textbf{Input:} The model's input consists of the user set 
% $\mathcal{U}$, the item set $\mathcal{V}$,  the observed interactions 
% ${\mathcal{V}}_{\mathcal{U}}$, 
%  the observed interactions' timestamp set $T = \{t_1, t_2, ..., t_n\}$ where 
%  $n$ 
% denotes the total number of timestamps, and the collaborative knowledge graph 
% $\mathcal{G}_R$.

\item \textbf{Output:} Given a user $u$ and a timestamp $\hat{T}$, a recommender model results in an item set $\hat{\mathcal{V}} \subseteq 
\mathcal{V} $ with the corresponding explainable reasoning paths $\tau_{u, 
\hat{\mathcal{V}}}$.  
In particular, for an item $\hat{v} \in \hat{\mathcal{V}}$, 
$\tau_{u, \hat{v}}$ is a multi-hop path, which connects user $u$ with item $\hat{v}$ to explain why $\hat{v}$ is suitable for $u$: 
$\tau_{u, \hat{v}}=[u 
\stackrel{r_{1}}{\longrightarrow} e_{1} \stackrel{r_{2}}{\longrightarrow} 
\ldots \stackrel{r_{k-1}}{\longrightarrow} e_{k-1} 
\stackrel{r_{k}}{\longrightarrow} \hat{v}]$, where $e_{i-1} 
\stackrel{r_{i}}{\longrightarrow} e_{i}$ represents $(e_{i-1}, r_i, e_i) \in 
\mathcal{G}_T, i \in [k]$.

\end{itemize} 

% 可能一些, 比如类间距离大类内距离小的假设应该放在这里.
% 还有那个多跳推理能够成功的假设.




